\section*{Hoofdstuk \ref{ch:b1:faults-folds-fabrics} --- Breuken, plooien en structuren}

\begin{solution}{exo:b1:faults-folds-fabrics:1}
(a) Een \omterm{def:b1:faults-folds-fabrics:fault-types}{normaalbreuk}. (b) Een \omterm{def:b1:faults-folds-fabrics:fault-types}{omgekeerde breuk}, bijna zuiver dip-slip. (c) Een linksdraaiende (sinistrale) \omterm{def:b1:faults-folds-fabrics:fault-types}{strike-slip-breuk}. (d) Een rechtsdraaiende (dextrale) \omterm{def:b1:faults-folds-fabrics:fault-types}{strike-slip-breuk}. (e) Een \omterm{def:b1:faults-folds-fabrics:fault-types}{schuine breuk}: de \omterm{def:b1:faults-folds-fabrics:fault-elements}{hangende wand} is omlaag en langs de \omterm{def:b1:geological-maps:attitude}{strijk} naar links bewogen; haar strike-slip-component, $\cos\ang{40} = 0.77$ van de \omterm{def:b1:faults-folds-fabrics:fault-elements}{netto verschuiving}, is groter dan haar dip-slip-component, $0.64$, dus heet zij een sinistrale breuk met een normale component. (f) Met een helling van \ang{25}, minder dan \ang{45}, is de \omterm{def:b1:faults-folds-fabrics:fault-types}{omgekeerde breuk} een \omterm{def:b1:faults-folds-fabrics:fault-types}{overschuiving}.
\end{solution}

\begin{solution}{exo:b1:faults-folds-fabrics:2}
Met $|\sin r| = 1$ is de \omterm{def:b1:cross-sections:throw}{worp} $300\sin\ang{55} = \qty{245.7}{m}$ en de
\omterm{def:b1:cross-sections:throw}{horizontale verplaatsing} $300\cos\ang{55} = \qty{172.1}{m}$.
\end{solution}

\begin{solution}{exo:b1:faults-folds-fabrics:3}
Voor een \omterm{def:b1:faults-folds-fabrics:fold-classes}{rechtopstaande plooi} waarvan de \omterm{def:b1:faults-folds-fabrics:fold-elements}{flanken} $\delta_1$ en $\delta_2$ in tegengestelde richtingen helen, geldt $\iota = \ang{180} - \delta_1 - \delta_2$. (a) \ang{110}: open. (b) \ang{35}: strak. (c) \ang{7}: isoklinaal.
\end{solution}

\begin{solution}{exo:b1:faults-folds-fabrics:4}
\omterm{def:b1:faults-folds-fabrics:fault-types}{Normaalbreuken} helen $\ang{45} + \ang{35}/2 = \ang{62.5}$, \omterm{def:b1:faults-folds-fabrics:fault-types}{overschuivingen}
$\ang{45} - \ang{17.5} = \ang{27.5}$. Geconjugeerde \omterm{def:b1:faults-folds-fabrics:fault-types}{strike-slip-breuken} maken elk
\ang{27.5} met $\sigma_1$, aan weerszijden ervan, zodat zij samenkomen onder \ang{55}.
\end{solution}

\begin{solution}{exo:b1:faults-folds-fabrics:5}
(a) De strike-slip-component van de breuk is $400\cos\ang{30} = \qty{346.4}{m}$ langs de \omterm{def:b1:geological-maps:attitude}{strijk}, de dip-slip-component $400\sin\ang{-30} = \qty{-200}{m}$, dat is
\qty{200}{m} omlaag langs de helling; de \omterm{def:b1:cross-sections:throw}{worp} is $200\sin\ang{50} = \qty{153.2}{m}$ en
de \omterm{def:b1:cross-sections:throw}{horizontale verplaatsing} $200\cos\ang{50} = \qty{128.6}{m}$. (b) De \omterm{def:b1:geological-maps:attitude}{strijk} is $000$, dus
$\mathbf s = (0, 1, 0)$; $\mathbf u = (-\cos\ang{50}, 0, \sin\ang{50}) =
(-0.643, 0, 0.766)$. Dan
$\overrightarrow{AA'} = 400(0.866\,\mathbf s - 0.5\,\mathbf u) =
(128.6, 346.4, -153.2)$ in meter. Haar trend is $\arctan(128.6/346.4) =
\ang{20.4}$ en haar plunge $\arcsin(153.2/400) = \ang{22.5}$, wat ook
$\arcsin(\sin\ang{30}\sin\ang{50})$ is. (c) Een \omterm{def:b1:faults-folds-fabrics:fault-types}{schuine breuk}: de \omterm{def:b1:faults-folds-fabrics:fault-elements}{rake} ligt
tussen $0$ en \ang{-90}, dichter bij $0$: een sinistrale breuk met een normale
component.
\end{solution}

\begin{solution}{exo:b1:faults-folds-fabrics:6}
(a) De twee breuken maken \ang{60} met elkaar; $\sigma_1$ halveert de scherpe hoek,
dus is zij horizontaal langs $050$ (en $230$); $\sigma_3$ is horizontaal
langs $140$; $\sigma_2$ is verticaal, langs de lijn waarin de twee verticale breuken
elkaar snijden. Elke breuk maakt \ang{30} met $\sigma_1$, dus $\ang{45} - \phi/2 =
\ang{30}$ en $\phi = \ang{30}$. (b) De twee breuken snijden de grond in vier
wiggen; de wiggen die $\sigma_1$ bevatten, naar $050$ en naar $230$, bewegen naar elkaar toe,
en die die $\sigma_3$ bevatten bewegen uiteen, naar $140$ en $320$. Neem de breuk die
$020$ loopt, tussen de wig naar $050$ (ten oosten ervan) en de wig naar $320$ (ten
westen ervan): ten opzichte van de westelijke wig beweegt de oostelijke naar $230$ en weg
van $320$, die beide een component naar $200$ langs de breuk hebben. Een waarnemer op de
westelijke wig, die naar het oosten kijkt, ziet de oostelijke naar $200$ bewegen, naar rechts:
de breuk die $020$ loopt is \omterm{def:b1:faults-folds-fabrics:fault-types}{dextraal}. Hetzelfde argument, of de symmetrie om
$\sigma_1$, maakt de breuk die $080$ loopt \omterm{def:b1:faults-folds-fabrics:fault-types}{sinistraal}.
\end{solution}

\begin{solution}{exo:b1:faults-folds-fabrics:7}
De opwaartse normalen zijn
\begin{align*}
  \mathbf n_1 &= (\sin\ang{40}\sin\ang{120}, \sin\ang{40}\cos\ang{120}, \cos\ang{40}) = (0.557, -0.321, 0.766),\\
  \mathbf n_2 &= (\sin\ang{30}\sin\ang{260}, \sin\ang{30}\cos\ang{260}, \cos\ang{30}) = (-0.492, -0.087, 0.866).
\end{align*}
Het uitwendig product van hun is
$(-0.211, -0.859, -0.207)$, dat omlaag wijst: de \omterm{def:b1:faults-folds-fabrics:fold-elements}{scharnierlijn} helt
$\arcsin 0.207/0.909 = \ang{13.1}$ naar $\ang{180} + \arctan(0.211/0.859) =
\ang{193.8}$. $\mathbf n_1\cdot\mathbf n_2 = 0.417$, de hoek tussen de
normalen is \ang{65.3}, en $\iota = \ang{114.7}$: een \omterm{def:b1:faults-folds-fabrics:fold-classes}{open plooi}.
\end{solution}

\begin{solution}{exo:b1:faults-folds-fabrics:8}
$\mathbf n_1 = (0.321, 0.383, 0.866)$ voor de gelaagdheid en
$\mathbf n_2 = (0.337, -0.925, 0.174)$ voor de \omterm{def:b1:faults-folds-fabrics:foliation}{klievage}. Hun uitwendig product,
omlaag genomen, geeft een lijn die \ang{25.4} helt naar \ang{074.8}. Volgens
\cref{prop:b1:faults-folds-fabrics:intersection} is dit de richting van de
\omterm{def:b1:faults-folds-fabrics:fold-elements}{scharnierlijn}: de \omterm{def:b1:geological-maps:fold}{plooi} helt ongeveer \ang{25} naar het oost-noordoosten, wat
elke andere \omterm{def:b1:geological-maps:geological-map}{ontsluiting} van een van beide \omterm{def:b1:faults-folds-fabrics:fold-elements}{flanken} moet bevestigen.
\end{solution}

\begin{solution}{exo:b1:faults-folds-fabrics:9}
(a) Een waarnemer ten zuiden van de breuk, die naar het noorden kijkt, ziet het noordelijke deel van de
gang naar het oosten verplaatst, naar rechts: de breuk is \omterm{def:b1:faults-folds-fabrics:fault-types}{dextraal}. (b) De gang is een
verticaal vlak loodrecht op de breuk, dus is haar verspringing langs de \omterm{def:b1:geological-maps:attitude}{strijk} de
strike-slip-component: \qty{350}{m}, wat de hele verschuiving is als de
\omterm{def:b1:faults-folds-fabrics:fault-rocks}{slickenlines} horizontaal zijn. (c) De verschuiving maakt dan \ang{20} met de \omterm{def:b1:geological-maps:attitude}{strijk}
in het verticale breukvlak: haar horizontale component is nog steeds \qty{350}{m},
dus is de \omterm{def:b1:faults-folds-fabrics:fault-elements}{netto verschuiving} $350/\cos\ang{20} = \qty{372.5}{m}$ en haar verticale
component $350\tan\ang{20} = \qty{127.4}{m}$. Het noordelijke blok is naar het oosten bewogen,
en de \omterm{def:b1:faults-folds-fabrics:fault-rocks}{slickenlines} dalen naar het oosten, dus is het ook omlaag gegaan, met
\qty{127.4}{m}. Een verticale gang kan die verticale beweging niet tonen, want als zij langs zichzelf wordt geschoven, blijft haar spoor onveranderd.
\end{solution}

\begin{solution}{exo:b1:faults-folds-fabrics:10}
(a) De normaal op de breuk maakt $\ang{90} - \theta$ met $\sigma_1$, dus volgens
\cref{prop:b1:stress-strain-rheology:mohr},
$\sigma_n' = \sigma_1'\sin^2\theta + \sigma_3'\cos^2\theta$ en
$\tau = (\sigma_1' - \sigma_3')\sin\theta\cos\theta$. De glijding begint wanneer
$\tau = \mu\sigma_n'$; delen door $\sigma_3'\cos^2\theta$ en schrijven
$R = \sigma_1'/\sigma_3'$, $(R - 1)\tan\theta = \mu(R\tan^2\theta + 1)$, dus
$R\tan\theta(1 - \mu\tan\theta) = \tan\theta + \mu$, wat de formule is.
(b) Met $\sigma_1$ verticaal is $\theta$ \ang{90} min de helling:
$\theta = \ang{30}$ geeft $R = 3.12$; $\theta = \ang{45}$, $R = 4.00$;
$\theta = \ang{55}$, $R = 9.92$. (c) De noemer verdwijnt wanneer
$\tan\theta = 1/\mu$, $\theta = \ang{59.0}$; daarbuiten zou $R$ negatief zijn:
geen verhouding van positieve \omterm{def:b1:stress-strain-rheology:lithostatic}{effectieve spanningen} kan de breuk doen glijden, hoe groot
$\sigma_1'$ ook wordt. \omterm{def:b1:faults-folds-fabrics:fault-types}{Normaalbreuken} die minder dan \ang{31.0} helen, kunnen
nooit glijden onder een verticale $\sigma_1$ met $\mu = 0.6$. (d) Het best georiënteerde vlak ligt op
$\ang{45} - \phi/2$ van $\sigma_1$. Voor $\mu = 0.6$, $\phi = \ang{31.0}$:
\omterm{def:b1:faults-folds-fabrics:fault-types}{normaalbreuken} helen \ang{60.5}, \omterm{def:b1:faults-folds-fabrics:fault-types}{overschuivingen} \ang{29.5}. Voor $\mu = 0.85$,
$\phi = \ang{40.4}$: \ang{65.2} en \ang{24.8}.
\end{solution}

\begin{solution}{exo:b1:faults-folds-fabrics:11}
Kijkend naar het noorden ligt het oosten rechts, en vlakken die naar het westen helen, stijgen naar
rechts. (a) Bij B helt de gelaagdheid op dezelfde manier als de \omterm{def:b1:faults-folds-fabrics:foliation}{klievage}, steiler:
volgens \cref{prop:b1:faults-folds-fabrics:cleavage-bedding}(a) ligt B op een
overgekantelde flank; A en C, waar de gelaagdheid minder steil is, staan rechtop.
(b) De \omterm{def:b1:faults-folds-fabrics:foliation}{klievage} stijgt onder \ang{60}; bij A en C stijgt de gelaagdheid onder
\ang{40}, en de scherpe rotatie van de een naar de ander is met de klok mee, over
\ang{20}, en verplaatst de top van de \omterm{def:b1:faults-folds-fabrics:foliation}{klievage} naar het oosten; bij B stijgt de gelaagdheid
onder \ang{80} en is de rotatie tegen de klok in, over \ang{20}, en verplaatst zij die naar
het westen. (c) De zin verandert tussen A en B en opnieuw tussen B en C: in elk interval
ligt een \omterm{def:b1:faults-folds-fabrics:fold-elements}{scharnier}. Bij A, dat rechtop staat, verplaatst de top van de
\omterm{def:b1:faults-folds-fabrics:foliation}{klievage} zich naar het oosten, naar de \omterm{def:b1:geological-maps:fold}{antiklinaal} aan de top van de flank: het \omterm{def:b1:faults-folds-fabrics:fold-elements}{scharnier}
tussen A en B is een \omterm{def:b1:geological-maps:fold}{antiklinaal}, en dat tussen B en C een \omterm{def:b1:geological-maps:fold}{synklinaal}.
(d) In de notatie van het bewijs helt de rechtopstaande \omterm{def:b1:faults-folds-fabrics:fold-elements}{flanken}
$\gamma - \iota/2 = \ang{40}$ en de overgekantelde \omterm{def:b1:faults-folds-fabrics:fold-elements}{flanken} $\gamma + \iota/2 =
\ang{80}$, dus $\gamma = \ang{60}$, zoals de \omterm{def:b1:faults-folds-fabrics:foliation}{klievage} zegt, en $\iota =
\ang{40}$: \omterm{def:b1:faults-folds-fabrics:fold-classes}{strakke plooien}. De axiale vlakken \omterm{def:b1:geological-maps:attitude}{strijken} noord en helen \ang{60} naar het
westen, ligging $270/60$. De korte \omterm{def:b1:faults-folds-fabrics:fold-elements}{flanken} zijn voorbij de verticale gedraaid met
hun toppen naar het oosten: de \omterm{def:b1:geological-maps:fold}{plooien} zijn naar het oosten vergent, zoals in de figuur van
het hoofdstuk.
\end{solution}

\begin{solution}{exo:b1:faults-folds-fabrics:12}
(a) In een symmetrische chevronplooi maakt elke flank de hoek $\iota/2$ met het axiale vlak, dat loodrecht staat op de richting van de verkorting; een flank van lengte $l$ overspant dus $l\sin(\iota/2)$ over de \omterm{def:b1:geological-maps:fold}{plooien}, en de hele laag $L_0\sin(\iota/2)$. (b) $\sin\ang{30} = 0.5$: een verkorting van
\qty{50}{\percent}. (c) $\sin(\iota/2) = 0.7$ geeft $\iota = \ang{88.9}$.
(d) $\iota = \ang{156}$ en $\sin\ang{78} = 0.978$: een verkorting van
\qty{2.2}{\percent}, wat ook $1 - \cos\ang{12}$ is.
\end{solution}

\begin{solution}{pb:b1:faults-folds-fabrics:1}
\textbf{1.} $\mathbf n = (\sin\delta\sin\alpha, \sin\delta\cos\alpha,
\cos\delta)$ geeft
\[
  \mathbf n_1 = (0.20855, 0.03415, 0.97742), \qquad
  \mathbf n_2 = (-0.20855, 0.03415, 0.97742).
\]

\textbf{2.} $\mathbf n_1\cdot\mathbf n_2 = -0.04349 + 0.00117 + 0.95535 =
0.91302$, en de hoek tussen de normalen is \ang{24.1}.

\textbf{3.} Volgens \cref{prop:b1:faults-folds-fabrics:hinge} geldt
$\iota = \ang{180} - \ang{24.1} = \ang{155.9}$.

\textbf{4.} Tussen \ang{170} en \ang{90}: een \omterm{def:b1:faults-folds-fabrics:fold-classes}{open plooi}.

\textbf{5.} De \omterm{def:b1:faults-folds-fabrics:fold-elements}{flanken} helen naar $080.7$ en $279.3$, niet precies oost en west,
omdat de \omterm{def:b1:geological-maps:fold}{plooi} dompelt: hun steilste lijnen liggen niet in het profiel, dat alleen de
\omterm{def:b1:cross-sections:apparent-dip}{schijnbare helling} toont. In een oost--west-doorsnede, volgens
\cref{thm:b1:cross-sections:apparent-dip}, is $\tan\beta = \tan\ang{12.2}
\cos\ang{9.3} = 0.2134$, dus $\beta = \ang{12.0}$, en $\ang{180} -
2\times\ang{12.0} = \ang{155.9}$ (het profiel is \ang{2} van die doorsnede gekanteld, wat bij deze nauwkeurigheid niets verandert).

\textbf{6.} Haar \omterm{def:b1:faults-folds-fabrics:fold-elements}{flanken} zijn evenwijdig aan die van de \omterm{def:b1:geological-maps:fold}{antiklinaal}, dus is haar \omterm{def:b1:faults-folds-fabrics:fold-elements}{interflankhoek} ook
\ang{155.9}.

\textbf{7.} $\mathbf n_1\times\mathbf n_2 = (0.03415\times0.97742 -
0.97742\times0.03415,\; 0.97742\times(-0.20855) - 0.20855\times0.97742,\;
0.20855\times0.03415 + 0.03415\times0.20855) = (0, -0.40768, 0.01424)$.
Omlaag genomen, $(0, 0.40768, -0.01424)$: de \omterm{def:b1:faults-folds-fabrics:fold-elements}{scharnierlijn} heeft trend $000$ en helt
$\arctan(0.01424/0.40768) = \ang{2.0}$ naar het noorden.

\textbf{8.} Het axiale vlak bevat de \omterm{def:b1:faults-folds-fabrics:fold-elements}{scharnierlijn} en halveert de twee \omterm{def:b1:faults-folds-fabrics:fold-elements}{flanken};
zijn normaal ligt langs $\mathbf n_1 - \mathbf n_2 = (0.41710, 0, 0)$, horizontaal
en oost--west. Het is verticaal en strijkt noord--zuid: $090/90$. De
\omterm{def:b1:geological-maps:fold}{antiklinaal} is rechtopstaand.

\textbf{9.} Een open, \omterm{def:b1:faults-folds-fabrics:fold-classes}{rechtopstaande plooi}, dompelend \ang{2.0} naar het noorden.

\textbf{10.} Volgens \cref{prop:b1:geological-maps:folds} sluiten de \omterm{def:b1:geological-maps:geological-map}{ontsluitingen} van een
\omterm{def:b1:geological-maps:fold}{antiklinaal} in de richting van de dompeling: naar het noorden. Haar kern bevat de
oudste gesteenten: het zandsteen aan de basis van het dek en, waar de rivier het diepst
heeft ingesneden, het fundament van schist en graniet daaronder.

\textbf{11.} Axiaalvlakkig staat de \omterm{def:b1:faults-folds-fabrics:foliation}{klievage} verticaal en strijkt zij noord--zuid,
$090/90$. Volgens \cref{prop:b1:faults-folds-fabrics:intersection} is haar \omterm{def:b1:faults-folds-fabrics:lineation}{snijlineatie}
met de gelaagdheid evenwijdig aan de \omterm{def:b1:faults-folds-fabrics:fold-elements}{scharnierlijn}: trend
$000$, plunge \ang{2.0}. Rechtstreeks: het vlak $080.7/12.2$ snijdt het verticale vlak dat noord strijkt langs een lijn met plunge $\arctan(\tan\ang{12.2}
\cos\ang{80.7}) = \ang{2.0}$ naar $000$.

\textbf{12.} Kijkend naar het noorden stijgt de oostflank naar links onder \ang{12.0}
en is de \omterm{def:b1:faults-folds-fabrics:foliation}{klievage} verticaal: de scherpe rotatie van de \omterm{def:b1:faults-folds-fabrics:foliation}{klievage} naar de
gelaagdheid, over \ang{78.0}, is tegen de klok in en verplaatst de top van de
\omterm{def:b1:faults-folds-fabrics:foliation}{klievage} naar het westen. De flank staat rechtop, dus volgens
\cref{prop:b1:faults-folds-fabrics:cleavage-bedding}(b) ligt het antiklinaalscharnier
naar het westen, bij $x = \qty{3.5}{km}$, zoals de kaart toont.

\textbf{13.} De \omterm{def:b1:geological-maps:attitude}{strijk} is $120 - 90 = 030$; het vlak helt naar $120$, dus
ligt de \omterm{def:b1:faults-folds-fabrics:fault-elements}{hangende wand} aan de oost-zuidoost-kant van het spoor.

\textbf{14.} De breuk is normaal: $\sigma_1$ was verticaal.

\textbf{15.} $\sigma_2$ is horizontaal langs de \omterm{def:b1:geological-maps:attitude}{strijk}, $030$--$210$;
$\sigma_3$ horizontaal dwars erop, $120$--$300$.

\textbf{16.} $\ang{45} + \phi/2 = \ang{60}$ geeft $\phi = \ang{30}$, en
$\mu = \tan\ang{30} = 0.577$.

\textbf{17.} De geconjugeerde set strijkt eveneens $030$ en helt \ang{60} de andere kant op: $300/60$. De twee sets komen samen onder $\ang{180} - 2\times\ang{60} = \ang{60}$,
met de verticale $\sigma_1$ die de scherpe hoek halveert.

\textbf{18.} Met $\mu = 0.85$, $\phi = \ang{40.4}$ en de best georiënteerde
\omterm{def:b1:faults-folds-fabrics:fault-types}{normaalbreuken} helen \ang{65.2}; de Varden-breuk, op \ang{60}, staat \ang{5.2} ernaast.
Volgens \cref{exo:b1:faults-folds-fabrics:10} glijdt zij bij de verhouding
$\sigma_1'/\sigma_3' = (1 + 0.85\cot\ang{30})/(1 - 0.85\tan\ang{30}) = 4.85$,
tegen $4.68$ voor het beste vlak: zij is bijna, maar niet precies, het best
georiënteerde, en zou opnieuw glijden bij een spanningsverhouding slechts \qty{4}{\percent}
hoger.

\textbf{19.} De verschuiving is omlaag langs de helling: $r = \ang{-90}$; de \omterm{def:b1:faults-folds-fabrics:fault-elements}{verschuivingsvector} heeft trend
$120$ en helt \ang{60}.

\textbf{20.} $(\cos\ang{60}\sin\ang{120}, \cos\ang{60}\cos\ang{120},
-\sin\ang{60}) = (0.4330, -0.2500, -0.8660)$.

\textbf{21.} $e = 200/\tan\ang{60} = \qty{115.5}{m}$.

\textbf{22.} Oost: $115.47\sin\ang{120} = \qty{100.0}{m}$; noord:
$115.47\cos\ang{120} = \qty{-57.7}{m}$, dat is \qty{57.7}{m} naar het zuiden.

\textbf{23.} $L = 200/\sin\ang{60} = \qty{230.9}{m}$, en
$200^2 + 115.47^2 = 40\,000 + 13\,333 = 53\,333 = 230.94^2$. De \omterm{def:b1:cross-sections:throw}{worp} en de
\omterm{def:b1:cross-sections:throw}{horizontale verplaatsing} zijn de verticale en horizontale componenten van de verschuiving, die in het verticale vlak van de helling ligt: zij zijn de zijden van een rechthoekige driehoek waarvan de schuine zijde de verschuiving is.

\textbf{24.} Op de meter nauwkeurig is de \omterm{def:b1:faults-folds-fabrics:fault-elements}{netto verschuiving} van de breuk van de Varden Hills
\textbf{\qty[text-series-to-math=true]{231}{m}}.
\end{solution}
